% Source-derived chapter 06: Abelian varieties over finite fields
% Original source: standard-conjectures-abelian-fourfolds
\chapter{Abelian varieties over finite fields}
\label{standard-conjectures-abelian-fourfolds-chapter-06}
\source{standard-conjectures-abelian-fourfolds}

\begin{remark}[Source section]
\label{standard-conjectures-abelian-fourfolds-chapter-06-source}
\source{standard-conjectures-abelian-fourfolds}
\source{standard-conjectures-abelian-fourfolds}
This chapter reproduces the corresponding section of the retrieved
LaTeX source, including its definitions, statements, examples, and
proofs. It has not yet been translated into Lean.
\notready
\end{remark}

% The source section title was: Abelian varieties over finite fields
\label{finitefields}
\source{standard-conjectures-abelian-fourfolds}




We start by recalling classical results on abelian varieties over finite fields and afterwards we draw some consequences on motives and algebraic cycles. The abundance of endomorphisms (due to Tate) will allow to decompose the motive of such an abelian variety in small factors. The main results  are Proposition \ref{talvez} and Corollary \ref{lifting motives} which will be essential to apply Theorem \ref{mainthm} to abelian varieties. The latter  says that such a decomposition lifts to characteristic zero and the former that each factor of the decomposition either does not contain algebraic classes (hence it is not interesting for our problem) either it is spanned by algebraic classes.

Throughout the section, we consider an abelian variety $A$ of dimension $g$ over a finite field $k$. We denote by $\End(A)$ the ring of endomorphisms of $A$ and by $\Frob_A$ the Frobenius of $A$. We write $\End^0(A)$ for $\End(A)\otimes_\Z \Q$,  and $*$ for the Rosati involution on it (induced by a polarization).

\begin{theorem}
\source{standard-conjectures-abelian-fourfolds}
\notready(Tate)\label{tate} A maximal  commutative $\Q$-subalgebra of $\End^0(A)$ has dimension $2g$ \cite{Tate}. There exist maximal commutative $\Q$-subalgebras of $\End^0(A)$ which are the product of CM fields \cite[Lemma 2]{TateBour}.\end{theorem}
\source{standard-conjectures-abelian-fourfolds}
\begin{remark}
\source{standard-conjectures-abelian-fourfolds}
\notready\label{remfrob}
\source{standard-conjectures-abelian-fourfolds}
Any maximal commutative $\Q$-subalgebra of $\End^0(A)$ must contain  $\Frob_A$ as the latter is contained in the center of $\End^0(A)$. 
\end{remark}
\begin{definition}
\source{standard-conjectures-abelian-fourfolds}
\notready\label{definitionCM} 
\source{standard-conjectures-abelian-fourfolds}
When $A$ is a simple abelian variety over $k$ a CM-structure for $A$ is the choice of a maximal CM field in $\End^0(A)$. 

For a general $A$ the choice of a decomposition of $A$  in the isogeny category as a product of simple abelian varieties $A_1\times \cdots\times A_t$ and a  CM-structure $L_i$ for each $A_i$ induce a maximal commutative $\Q$-subalgebra ${B=L_1\times \cdots\times L_t}$ of $\End^0(A)$. An algebra $B$  constructed as above  is called a CM-structure for $A$. If such an algebra is $*$-stable we will say that the CM-structure is $*$-stable.
\end{definition}
\begin{remark}
\source{standard-conjectures-abelian-fourfolds}
\notready
We will see (Corollary \ref{CMstab}) that given a CM-structure there exists (possibly after a finite extension of the base field $k$ and after isogeny) a polarization for which  the CM-structure is $*$-stable.
\end{remark}


\begin{notation}
\source{standard-conjectures-abelian-fourfolds}
\notready\label{NotationCM}
\source{standard-conjectures-abelian-fourfolds}
Let $B$ be a CM-structure for $A$. We write \[{B=L_1\times \cdots\times L_t}\] as a product of CM fields. Let 
 $L$ be the number field which is the Galois closure (over $\Q$) of the compositum of the fields $L_i$, it is a CM number field as well \cite[Proposition 5.12]{goro}. Let $\Sigma_i$ be the set of morphisms from $L_i$ to $L$ and   $\Sigma$ be the disjoint union 
 \[\Sigma = \Sigma_1\cup\cdots \cup \Sigma_t.\]
Write $\bar{\cdot}$ for the action on $\Sigma$ induced by composition with  complex conjugation. We will use the same notation for the induced action on subsets of  $\Sigma$.
\end{notation}


\begin{prop}
\source{standard-conjectures-abelian-fourfolds}
\notready\label{decomposition}
\source{standard-conjectures-abelian-fourfolds}
Let $L$ and   $\Sigma$ be as in Notation  \ref{NotationCM}. Then, in $\CHM (k)_{L}$, the motive $\mathfrak{h}^1(A)$ decomposes into a sum of $2g$ motives  
\[\mathfrak{h}^1(A) =\bigoplus_{\sigma \in \Sigma} M_\sigma,\]
where the action of $b\in L_i$ on $M_\sigma$ induced by Theorem \ref{classic}(\ref{kings}) is given by multiplication by $\sigma(b)$ if $\sigma \in \Sigma_i$ and by multiplication by zero otherwise.  Moreover, each motive $ M_\sigma$ is of rank one (in the sense of Definition \ref{defrank}). Finally, if the CM-structure is $*$-stable, the isomorphism $p:  \mathfrak{h}^1(A)\cong\mathfrak{h}^1(A)^{\vee}(-1)$ of Theorem \ref{classic}(\ref{polarization}) restricts to an isomorphism \[M_\sigma\cong M_{\bar{\sigma}}^\vee(-1)\] for all $\sigma\in \Sigma$, and to the zero map 
\[ M_\sigma \stackrel{0}{\longrightarrow} M_{\sigma'}^\vee(-1) \]
for all $\sigma'\neq {\bar{\sigma}}$.
\end{prop}
\begin{proof}
This is \cite[Corollary 3.2]{Ascona}, we recall here the argument. Consider the inclusion $L_1\times \cdots\times L_t  \hookrightarrow  \End^0(A).$ By Theorem \ref{classic}(\ref{kings}), we deduce an inclusion $(\prod_i L_i)\otimes L  \hookrightarrow  \End_{\CHM^{\ab}(k)_{L}}(\mathfrak{h}^1(A)).$ Each projector of the product of fields $(\prod_i L_i)\otimes L \cong \prod_i L^{\Sigma_i}$ defines a factor $M_{\sigma}.$

Let us now consider the last part of the statement. First, it can be checked after realization   \cite[Corollary 1.12]{Ascona}. Secondly, Rosati involution acts as  complex conjugation on CM fields  \cite[pp. 211-212]{MumfordTata}. Then the statement follows by the very definition of Rosati involution.\end{proof}

%\begin{remark}
%In order to prove the Standard Conjecture of Hodge type we can work with $\langle \cdot, \cdot \rangle_{1,\mot}^{\otimes 2n}$ instead of $\langle \cdot, \cdot \rangle_{2n,\mot}$ thanks to Lemma \ref{accoppiamenti1}. This explains the interest of the following results, they will allow to work with  motives that are smaller than the motive of $A$.
%\end{remark}

\begin{proposition}
\source{standard-conjectures-abelian-fourfolds}
\notready\label{primadecomposizione}
\source{standard-conjectures-abelian-fourfolds}
In the setting of Notation \ref{NotationCM} and Proposition \ref{decomposition}  the following holds:

\begin{enumerate}


\item\label{decomposizioneL} In $\CHM (k)_{L}$ the motive $\mathfrak{h}^n(A)$ decomposes into a sum 
\source{standard-conjectures-abelian-fourfolds}
\[\mathfrak{h}^n(A) =\bigoplus_{I\subset \Sigma, \vert I \vert = n} M_I,\]
with $M_I= \otimes_{\sigma \in I}M_\sigma$. Each motive $M_I$ is of rank one (in the sense of Definition \ref{defrank}).

Moreover, if the CM-structure is $*$-stable, then the motives $M_I$ and $M_J$ are mutually orthogonal  in $\NUM(k)_L$ with respect to $\langle \cdot, \cdot \rangle_{1,\mot}^{\otimes n}$ (Definition \ref{defpairingmot}) except if $I = \overline{J}$.

\item\label{decomposizioneQ} In $\CHM (k)_{\Q}$ the motive $\mathfrak{h}^n(A)$ decomposes into a sum of motives
\source{standard-conjectures-abelian-fourfolds}
\[\mathfrak{h}^n(A) =\bigoplus  \mathcal{M}_I,\]
where $\mathcal{M}_I$ is spanned by the factors of the form $M_{g(I)}$, with $g$ varying in $\Gal(L/\Q).$

Moreover, if the CM-structure is \nobreak{$*$-stable}, this decomposition in $\NUM(k)_\Q$  is orthogonal with respect to $\langle \cdot, \cdot \rangle_{1,\mot}^{\otimes n}$  (Definition \ref{defpairingmot}).

\end{enumerate}
\end{proposition}
\begin{proof}
Part (\ref{decomposizioneQ}) follows from part (\ref{decomposizioneL}). For the latter, Proposition \ref{decomposition} gives the case $n=1$. Using Theorem \ref{classic}(\ref{kunn}) we deduce the result for higher $n$.
\end{proof}

\begin{proposition}
\source{standard-conjectures-abelian-fourfolds}
\notready\label{talvez}
\source{standard-conjectures-abelian-fourfolds}
Let  $\mathcal{M}_I \in \CHM (k)_{\Q}$ be a  direct factor of $\mathfrak{h}^{2n}(A)$ as constructed in Proposition \ref{primadecomposizione}(\ref{decomposizioneQ}). Then the following holds:
\begin{enumerate}
\item\label{generatoalgebrico} If the vector space $\Hom_{\NUM(k)_\Q}(\one, \mathcal{M}_I (n))$ is not zero then 
\source{standard-conjectures-abelian-fourfolds}
\[\dim_\Q \Hom_{\NUM(k)_\Q}(\one, \mathcal{M}_I (n)) = \dim \mathcal{M}_I.\]
In this case homological and numerical equivalence coincide on $ \mathcal{M}_I$ and the realizations of $ \mathcal{M}_I$ are spanned by algebraic classes.
\item\label{generatolefschetz}   If the vector space $\Hom_{\NUM(k)_\Q}(\one, \mathcal{M}_I (n))$ contains a nonzero Lefschetz class (Definition \ref{defexotic}), then all algebraic  classes in 
\source{standard-conjectures-abelian-fourfolds}
it are Lefschetz.
\end{enumerate}
\end{proposition}
\begin{proof}
As numerical equivalence commutes with  extension of scalars over $\mathbb{Q}$ \cite[Proposition 3.2.7.1]{Andmot}, we have \[\dim_\Q \Hom_{\NUM(k)_\Q}(\one, \mathcal{M}_I (n)) = \dim_L \Hom_{\NUM(k)_L}(\one, \mathcal{M}_I (n)).\]
 On the other hand, by construction of $\mathcal{M}_I$ (see Proposition \ref{primadecomposizione}(\ref{decomposizioneQ})), we have that the space $\Hom_{\NUM(k)_L}(\one,\mathcal{M}_I(n))$ is generated by the spaces $\Hom_{\NUM(k)_L}(\one,M_{g(I)}(n))$, with $g$ varying in $\Gal(L/\Q).$

 Moreover, as $\Gal(L/\Q)$  acts on the space of algebraic cycles modulo numerical equivalence, we have 
 \[\dim_L \Hom_{\NUM(k)_L}(\one,M_{I}(n)) = \dim_L \Hom_{\NUM(k)_L}(\one,M_{g(I)}(n)),\] for all $g\in \Gal(L/\Q).$
 
 Now, by Proposition \ref{primadecomposizione}(\ref{decomposizioneL}), the motives $M_{I} \in \CHM (k)_{L}$ are of rank one hence, by Proposition \ref{motivisupersingolari}, the above  dimension is either zero or one. This implies the equality in part (\ref{generatoalgebrico}). The rest of part (\ref{generatoalgebrico}) follows from Proposition \ref{motivisupersingolari} as well. 
 
 
 For part (\ref{generatolefschetz}), the proof just goes as part (\ref{generatoalgebrico}) as all the properties of motives and algebraic classes that we used hold true for Lefschetz motives and Lefschetz classes by   \cite[p. 640]{Milnelef}.
\end{proof}
\begin{remark}
\source{standard-conjectures-abelian-fourfolds}
\notready The previous proposition is not accessible if one replaces numerical with homological equivalence as it is not known that homological equivalence commutes with extension of scalars. One finds the same issues in \cite{Clozel}. This is the crucial reason for which the main results in this paper are in the setting of numerical equivalence.
\end{remark}
\begin{theorem}[{\cite[Theorem 2]{TateBour}}]
\source{standard-conjectures-abelian-fourfolds}
\notready\label{HondaTate}
\source{standard-conjectures-abelian-fourfolds}
 For any CM-structure $B$ for $A$ (Definition \ref{definitionCM}), the pair  $(A,B)$ lifts to characteristic zero. More precisely there exists a $p$-adic field $K$ with ring of integers $W$ whose residue field $k'$ is a finite extension of $k$, and there is an abelian scheme $\mathcal{A}$ over $W$ together with an embedding $B\hookrightarrow \End^0(\mathcal{A})$, such that $\mathcal{A}\times_W k'$ is isogenous to $A \times_k k'$ (and the isogeny is $B$-equivariant).
 \end{theorem}

\begin{corollary}
\source{standard-conjectures-abelian-fourfolds}
\notready\label{CMstab}
\source{standard-conjectures-abelian-fourfolds}
There exists a polarization on the abelian scheme $\mathcal{A}$ for which the CM-structure is $*$-stable
\end{corollary}
\begin{proof}
Consider the decomposition $\mathcal{A}= \mathcal{A}_1\times \cdots \times \mathcal{A}_t  $ in product  CM simple abelian schemes given from the choice of the CM-structure (Definition \ref{definitionCM}). Choose the polarization to be the product of a polarization on each factor. The statement then reduces to the case when $\mathcal{A}$ is simple and it is enough to check it on the generic fiber, hence in characteristic zero. On the other hand the endomorphism algebra of a simple CM abelian variety in characteristic zero cannot be bigger than the CM-structure itself \cite[\S 22]{MumfordTata}, hence the $*$-stability is automatic. 
\end{proof}


\begin{corollary}
\source{standard-conjectures-abelian-fourfolds}
\notready\label{lifting motives}
\source{standard-conjectures-abelian-fourfolds}
The decompositions of $\mathfrak{h}(A \times_k k')$ in Proposition \ref{primadecomposizione}  lift to decompositions of $\mathfrak{h}(\mathcal{A})\in  \CHM (W)$. Moreover, if the CM-structure of $A$ is $*$-stable (and if the polarisation lifts as well) then the orthogonality statement in Proposition \ref{primadecomposizione} holds true in $ \CHM (W)$ as well.
%More precisely, there is a $p$-adic field $K$ with ring of integers $W$ whose residue field, $k'$, is a finite extension of $k$, and there is a motive $ \tilde{\mathcal{M}_I} \in \CHM(W)_\Q$ such that the restriction  $ (\tilde{\mathcal{M}_I} )_{\vert{k'}} \in \CHM(k')_\Q$ is isomorphic to the base change $\mathcal{M}_I \times_k k'.$
\end{corollary}

%For any choice of $B$ as above the pair  $(A,B)$ lifts to characteristic zero. More precisely there exists a $p$-adic field $K$ with ring of integers $W$ whose residue field $k'$ is a finite extension of $k$, and there is an abelian scheme $\mathcal{A}$ over $W$ such that $B=\End^0(\mathcal{A})$ and $\mathcal{A}\times_W k$ is isogenous to $A \times_k k'.$


\begin{proof}
The proof of Proposition \ref{primadecomposizione}    is a formal combination of Theorem \ref{classic} together with   the CM-structure.
It works also over $W$ because of Remark \ref{AHP} and Theorem \ref{HondaTate}.
\end{proof}
